\section{Data Generation}

\renewcommand{\thefigure}{B.\arabic{figure}}  % Number figures as A.1, A.2, ...
\renewcommand{\thetable}{B.\arabic{table}}  % Number tables as A.1, A.2, ...


For a PDE defined over a spatial domain \( \mathcal{X} = \mathbb{R}^2 \) with \( \mathbf{x} = (x_1, x_2) \in \mathcal{X} \) and time \( t \in [0, T] \), including \( n_t \) discrete time steps of the solution over the interval \( [0, T] \), we generate training data for the Sensitivity-Constrained Neural Operator (SC-NO) to learn the mapping \( \mathcal{G}_\theta: \mathcal{U} \to \mathcal{V} \), where \(\mathcal{U}\) is the space of input functions and \(\mathcal{V}\) is the space of solutions, as described in the Methods section. Each data sample comprises an input function \( a(\textbf{x}, t, p) \in \mathcal{U} \), the corresponding solution \( u(\textbf{x}, t, p) = \mathcal{G}(a)(\textbf{x}, t, p) \in \mathcal{V} \), and the true Jacobian \( \frac{\partial u}{\partial p} \), computed from a high-fidelity solver, to support sensitivity-aware training.



\subsubsection*{Structure of the Input Function}

The input function is defined as:
\begin{equation}
a(\textbf{x}, t, p) = \big( a_{\text{init}}(\textbf{x}, t), p(\textbf{x}) \big),
\end{equation}
where:
\begin{itemize}
    \item \( a_{\text{init}}(\textbf{x}, t) \) contains the initial time steps of the solution:
    \begin{equation}
    a_{\text{init}}(\textbf{x}, t) = \big( u(\textbf{x}, t_1, p), u(\textbf{x}, t_2, p), \dots, u(\textbf{x}, t_{t_u}, p) \big), \quad t \in [0, t_u],
    \end{equation}
    representing \( t_u \) steps of the state \( u(\textbf{x}, t, p) \).


    \item \( p(\textbf{x}) \) is the spatially varying parameter field.
    

\end{itemize}
The total simulation time is:
\begin{equation}
n_t = t_u + t_G,
\end{equation}
where \( t_u \) is the number of initial time steps provided in \( a_{\text{init}} \), and \( t_G \) is the additional time steps predicted by \( \mathcal{G}_\theta \).

\subsubsection*{Dataset Structure with True Sensitivities}

Each dataset sample is structured as:
\begin{equation}
\big[ a(\textbf{x}, t, p),\; \textbf{x},\; \mathcal{G}(a)(\textbf{x}, t, p),\; \frac{\partial \mathcal{G}(a)}{\partial p}(\textbf{x}, t, p) \big],
\end{equation}
where:
\begin{itemize}
    \item \( a(\textbf{x}, t, p) = (a_{\text{init}}(x, t), p(\textbf{x})) \) is the input function in \(\mathcal{U}\).
    \item \( \textbf{x} = (x_1, x_2), \dots, (x_{n_x}, x_{n_y}) \) are spatial grid points over \(\mathcal{X}\).
    \item \( \mathcal{G}(a)(\textbf{x}, t, p) = u(\textbf{x}, t, p) \) is the true PDE solution over \( t \in [t_u, n_t] \).
    \item \( \frac{\partial \mathcal{G}(a)}{\partial p}(\textbf{x}, t, p) \) is the true Jacobian with respect to \( p(\textbf{x}) \), obtained from a high-fidelity solver.
\end{itemize}

For each realization \( i \), the input is:
\begin{equation}
a^{(i)}(\textbf{x}, t, p) =
\begin{bmatrix}
a_{\text{init}}^{(i)}(x_1, x_2, t) & a_{\text{init}}^{(i)}(x_2, x_3, t) & \cdots & a_{\text{init}}^{(i)}(x_{n_x}, x_{n_y}, t) \\
p^{(i)}(x_1, x_2) & p^{(i)}(x_2, x_3) & \cdots & p^{(i)}(x_{n_x}, x_{n_y})
\end{bmatrix}, \quad n_x \times n_y \times t_u,
\end{equation}
and the output solution is:
\begin{equation}
\mathcal{G}(a^{(i)})(\textbf{x}, t, p) =
\begin{bmatrix}
u^{(i)}(x_1, x_2, t, p(x_1, x_2)) \\
u^{(i)}(x_2, x_3, t, p(x_2, x_3)) \\
\vdots \\
u^{(i)}(x_{n_x}, x_{n_y}, t, p(x_{n_x}, x_{n_y}))
\end{bmatrix}, \quad n_x \times n_y \times t_G.
\end{equation}

Since \( u(\textbf{x}, t, p) \) evolves over space and time, its sensitivity to the spatially varying parameter \( p(\textbf{x}) \) is a matrix-valued function. To capture the cumulative effect of \( p(\textbf{x}) \) on the solution, we compute the Jacobian at the final time step \( t = T \):
\begin{equation}
\frac{\partial \mathcal{G}(a)}{\partial p}(\textbf{x}, T, p) =
\begin{bmatrix}
\frac{\partial u}{\partial p}(x_1, x_2; x_1, x_2) & \frac{\partial u}{\partial p}(x_1, x_2; x_2, x_3) & \cdots & \frac{\partial u}{\partial p}(x_1, x_2; x_{n_x}, x_{n_y}) \\[6pt]
\frac{\partial u}{\partial p}(x_2, x_3; x_1, x_2) & \frac{\partial u}{\partial p}(x_2, x_3; x_2, x_3) & \cdots & \frac{\partial u}{\partial p}(x_2, x_3; x_{n_x}, x_{n_y}) \\[6pt]
\vdots & \vdots & \ddots & \vdots \\[6pt]
\frac{\partial u}{\partial p}(x_{n_x}, x_{n_y}; x_1, x_2) & \frac{\partial u}{\partial p}(x_{n_x}, x_{n_y}; x_2, x_3) & \cdots & \frac{\partial u}{\partial p}(x_{n_x}, x_{n_y}; x_{n_x}, x_{n_y})
\end{bmatrix}, \quad (n_x \times n_y) \times (n_x \times n_y).
\end{equation}
Each entry \( \frac{\partial u}{\partial p}(x_i, x_j; x_k, x_l) \) quantifies the response of the solution at \( (x_i, x_j) \) at \( t = T \) to changes in \( p(x_k, x_l) \), reflecting both local and nonlocal PDE-induced dependencies. By evaluating the Jacobian at \( t = T \), we encapsulate the temporal propagation of sensitivity from \( t = 0 \) to \( T \), as the final state \( u(x, T, p) \) integrates the system’s evolution. This approach aligns with SC-NO training, where \( L_s = \ell_s\left( \frac{\partial \hat{u}}{\partial p}, \frac{\partial u}{\partial p} \right) \) enforces sensitivity accuracy, reducing computational cost by focusing on the final time step while preserving critical spatial-parametric relationships for optimization.

\subsection*{Scaled Gaussian Random Field for Parameter Generation}

To generate spatially varying random parameters \( p(\mathbf{x}) \) across the computational domain, we define a scaled and bounded Gaussian random field (GRF) as:

\begin{equation}
p(\mathbf{x}) = \mathcal{T} \big( Z(\mathbf{x}); \text{p}_{\text{min}}, \text{p}_{\text{max}}, \sigma, \mathcal{X} \big),
\end{equation}

where \( \mathcal{T} \) is a transformation ensuring that \( p(\mathbf{x}) \) is confined within the predefined range \( [\text{p}_{\text{min}}, \text{p}_{\text{max}}] \). Specifically, the transformation is defined as:

\begin{equation}
\mathcal{T} \big( Z(\mathbf{x}) \big) = \text{p}_{\text{min}} + (\text{p}_{\text{max}} - \text{p}_{\text{min}}) \cdot 
\frac{\tanh \left( \sigma Z(\mathbf{x}) \right) - \min\limits_{\mathbf{y} \in \mathcal{X}} \tanh \left( \sigma Z(\mathbf{y}) \right)}
{\max\limits_{\mathbf{y} \in \mathcal{X}} \tanh \left( \sigma Z(\mathbf{y}) \right) - \min\limits_{\mathbf{y} \in \mathcal{X}} \tanh \left( \sigma Z(\mathbf{y}) \right)}.
\end{equation}

Here, \( Z(\mathbf{x}) \) is a zero-mean Gaussian random field sampled as:

\begin{equation}
Z(\mathbf{x}) \sim \mathcal{N} \left( 0, (-\Delta + 9I)^{-2} \right),
\end{equation}

where \( \Delta \) is the Laplacian operator, enforcing spatial smoothness and local correlation in \( p(\mathbf{x}) \). The parameter \( \sigma \) is a scaling factor that modulates the range of fluctuations in the transformed field. The transformation \( \mathcal{T} \) normalizes the GRF realization using the hyperbolic tangent function \( \tanh(\cdot) \), ensuring smooth transitions and bounded outputs. The resulting field \( p(\mathbf{x}) \) maintains spatial coherence while adhering to the predefined range, making it well-suited for parameterizing complex PDE systems. By employing this approach, we achieve a structured yet stochastic representation of spatially varying parameters, facilitating robust generalization in neural operator training.
